\chapter{Hasil Eksperimen dan Analisis}

Bab ini menyajikan hasil pengujian setiap algoritma. Isi setiap penanda \textit{[...]} dengan hasil pengujian aktual dan rujuk grafik yang dihasilkan pada kolom terkait.

\renewcommand{\arraystretch}{1.15}
\newcommand{\hasil}[1]{\textit{[#1]}}
\newcommand{\awaltabel}{%
\begin{longtable}{|p{4.0cm}|p{10.2cm}|}
\hline
\textbf{Informasi yang dicatat} & \textbf{Data hasil pengujian} \\
\hline
\endfirsthead
\hline
\textbf{Informasi yang dicatat} & \textbf{Data hasil pengujian} \\
\hline
\endhead
\hline
\endfoot
\endlastfoot
}
\newcommand{\akhirtabel}{\end{longtable}}
\newcommand{\isiHC}{%
\textit{State} awal & \hasil{ringkasan/rujukan visual state awal} \\ \hline
\textit{State} akhir & \hasil{ringkasan/rujukan visual state akhir} \\ \hline
Nilai \textit{objective function} akhir & \hasil{nilai objektif akhir} \\ \hline
Grafik \textit{objective function} terhadap iterasi & \hasil{nomor gambar/rujukan grafik} \\ \hline
Durasi pencarian & \hasil{durasi (detik)} \\ \hline
Banyak iterasi hingga berhenti & \hasil{jumlah iterasi} \\ \hline
}
\newcommand{\isiGA}{%
\textit{State} awal (populasi awal) & \hasil{ringkasan/rujukan visual populasi awal} \\ \hline
\textit{State} akhir terbaik & \hasil{ringkasan/rujukan visual state akhir} \\ \hline
Nilai \textit{objective function} akhir terbaik & \hasil{nilai objektif akhir} \\ \hline
Grafik objektif maksimum dan rata-rata terhadap iterasi & \hasil{nomor gambar/rujukan grafik} \\ \hline
Jumlah populasi & \hasil{nilai parameter populasi} \\ \hline
Banyak iterasi & \hasil{nilai parameter iterasi} \\ \hline
Durasi pencarian & \hasil{durasi (detik)} \\ \hline
}

\newcommand{\percobaanHC}[2]{%
\subsubsection{Test case #1 --- Percobaan #2}
\awaltabel
\isiHC
\akhirtabel
}

% Argumen: nomor percobaan, awalan nama berkas gambar, nilai akhir, durasi, jumlah iterasi
\newcommand{\percobaanSteepest}[6]{%
\subsubsection{Test case #1 --- Percobaan #2}
\awaltabel
\textit{State} awal & \includegraphics[width=\linewidth]{gambar/pengujian/steepest-ascent-hc/#1/#3_initial_state_3d.png} \\ \hline
\textit{State} akhir & \includegraphics[width=\linewidth]{gambar/pengujian/steepest-ascent-hc/#1/#3_final_state_3d.png} \\ \hline
Nilai \textit{objective function} akhir & #4 \\ \hline
Grafik \textit{objective function} terhadap iterasi & \includegraphics[width=\linewidth]{gambar/pengujian/steepest-ascent-hc/#1/#3_experiment.png} \\ \hline
Durasi pencarian & #5 detik \\ \hline
Banyak iterasi hingga berhenti & #6 \\ \hline
\akhirtabel
}

\newcommand{\percobaanSideways}[2]{%
\subsubsection{Test case #1 --- Percobaan #2}
\awaltabel
\textit{Maximum sideways move} & \hasil{nilai parameter} \\ \hline
\textit{State} awal & \hasil{ringkasan/rujukan visual state awal} \\ \hline
\textit{State} akhir & \hasil{ringkasan/rujukan visual state akhir} \\ \hline
Nilai \textit{objective function} akhir & \hasil{nilai objektif akhir} \\ \hline
Grafik \textit{objective function} terhadap iterasi & \hasil{nomor gambar/rujukan grafik} \\ \hline
Durasi pencarian & \hasil{durasi (detik)} \\ \hline
Banyak iterasi hingga berhenti & \hasil{jumlah iterasi} \\ \hline
\akhirtabel
}
% Argumen: test case, nomor percobaan, awalan gambar, nilai akhir, durasi,
% iterasi, dan jumlah sideways move.
\newcommand{\percobaanSidewaysGambar}[7]{%
\subsubsection{Test case #1 --- Percobaan #2}
\awaltabel
\textit{Maximum sideways move} & 10 \\ \hline
\textit{State} awal & \includegraphics[width=\linewidth]{gambar/pengujian/sideways-move-hc/#1/#3_initial_state_3d.png} \\ \hline
\textit{State} akhir & \includegraphics[width=\linewidth]{gambar/pengujian/sideways-move-hc/#1/#3_final_state_3d.png} \\ \hline
Nilai \textit{objective function} akhir & #4 \\ \hline
Grafik \textit{objective function} terhadap iterasi & \includegraphics[width=\linewidth]{gambar/pengujian/sideways-move-hc/#1/#3_experiment.png} \\ \hline
Durasi pencarian & #5 detik \\ \hline
Banyak iterasi hingga berhenti & #6 \\ \hline
Jumlah \textit{sideways move} yang dilakukan & #7 \\ \hline
\akhirtabel
}
\newcommand{\percobaanRestart}[2]{%
\subsubsection{Test case #1 --- Percobaan #2}
\awaltabel
\textit{Maximum restart} & \hasil{nilai parameter} \\ \hline
\textit{State} awal & \hasil{ringkasan/rujukan visual state awal} \\ \hline
\textit{State} akhir terbaik & \hasil{ringkasan/rujukan visual state akhir} \\ \hline
Nilai \textit{objective function} akhir & \hasil{nilai objektif akhir} \\ \hline
Grafik \textit{objective function} terhadap iterasi & \hasil{nomor gambar/rujukan grafik} \\ \hline
Durasi pencarian & \hasil{durasi (detik)} \\ \hline
Banyak \textit{restart} yang dilakukan & \hasil{jumlah restart} \\ \hline
Banyak iterasi per \textit{restart} & \hasil{daftar/rata-rata iterasi per restart} \\ \hline
\akhirtabel
}
% Argumen: test case, nomor percobaan, awalan gambar, nilai akhir, durasi,
% banyak restart, iterasi per restart.
\newcommand{\percobaanRestartGambar}[7]{%
\subsubsection{Test case #1 --- Percobaan #2}
\awaltabel
\textit{Maximum restart} & 3 \\ \hline
\textit{State} awal & \includegraphics[width=\linewidth]{gambar/pengujian/random-restart-hc/#1/#3_initial_state_3d.png} \\ \hline
\textit{State} akhir terbaik & \includegraphics[width=\linewidth]{gambar/pengujian/random-restart-hc/#1/#3_final_state_3d.png} \\ \hline
Nilai \textit{objective function} akhir & #4 \\ \hline
Grafik \textit{objective function} terhadap iterasi & \includegraphics[width=\linewidth]{gambar/pengujian/random-restart-hc/#1/#3_experiment.png} \\ \hline
Durasi pencarian & #5 detik \\ \hline
Banyak \textit{restart} yang dilakukan & #6 \\ \hline
Banyak iterasi per \textit{restart} & #7 \\ \hline
\akhirtabel
}
% Argumen: test case, nomor percobaan, awalan gambar, nilai akhir, durasi, banyak iterasi
\newcommand{\percobaanStochasticGambar}[6]{%
\subsubsection{Test case #1 --- Percobaan #2}
\awaltabel
\textit{State} awal & \includegraphics[width=\linewidth]{gambar/pengujian/stochastic-hc/#1/#3_initial_state_3d.png} \\ \hline
\textit{State} akhir & \includegraphics[width=\linewidth]{gambar/pengujian/stochastic-hc/#1/#3_final_state_3d.png} \\ \hline
Nilai \textit{objective function} akhir & #4 \\ \hline
Grafik \textit{objective function} terhadap iterasi & \includegraphics[width=\linewidth]{gambar/pengujian/stochastic-hc/#1/#3_experiment.png} \\ \hline
Durasi pencarian & #5 detik \\ \hline
Banyak iterasi hingga berhenti & #6 \\ \hline
\akhirtabel
}
\newcommand{\percobaanSA}[2]{%
\subsubsection{Test case #1 --- Percobaan #2}
\awaltabel
\textit{State} awal & \hasil{ringkasan/rujukan visual state awal} \\ \hline
\textit{State} akhir & \hasil{ringkasan/rujukan visual state akhir} \\ \hline
Nilai \textit{objective function} akhir & \hasil{nilai objektif akhir} \\ \hline
Grafik \textit{objective function} terhadap iterasi & \hasil{nomor gambar/rujukan grafik} \\ \hline
Grafik $e^{\Delta E/T}$ terhadap iterasi & \hasil{nomor gambar/rujukan grafik} \\ \hline
Frekuensi \textit{stuck} pada \textit{local optima} & \hasil{jumlah/frekuensi kejadian} \\ \hline
Durasi pencarian & \hasil{durasi (detik)} \\ \hline
\akhirtabel
}
% Argumen: test case, nomor percobaan, awalan gambar, nilai akhir, frekuensi stuck, durasi
\newcommand{\percobaanSAGambar}[6]{%
\subsubsection{Test case #1 --- Percobaan #2}
\awaltabel
\textit{State} awal & \includegraphics[width=\linewidth]{gambar/pengujian/sa/#1/#3_initial_state_3d.png} \\ \hline
\textit{State} akhir & \includegraphics[width=\linewidth]{gambar/pengujian/sa/#1/#3_final_state_3d.png} \\ \hline
Nilai \textit{objective function} akhir & #4 \\ \hline
Grafik \textit{objective function} terhadap iterasi & \includegraphics[width=\linewidth]{gambar/pengujian/sa/#1/#3_experiment.png} \\ \hline
Grafik $e^{\Delta E/T}$ terhadap iterasi & Tercakup pada panel ke-4 grafik multi-panel di atas \\ \hline
Frekuensi \textit{stuck} pada \textit{local optima} & #5 episode \\ \hline
Durasi pencarian & #6 detik \\ \hline
\akhirtabel
}
\newcommand{\percobaanGA}[3]{%
\subsubsection{Test case #1 --- #2 --- Percobaan #3}
\awaltabel
\isiGA
\akhirtabel
}
% Argumen: test case, label konfigurasi, nomor percobaan, nilai akhir, populasi, iterasi, durasi
\newcommand{\percobaanGANilai}[7]{%
\subsubsection{Test case #1 --- #2 --- Percobaan #3}
\awaltabel
\textit{State} awal (populasi awal) & \hasil{ringkasan/rujukan visual populasi awal} \\ \hline
\textit{State} akhir terbaik & \hasil{ringkasan/rujukan visual state akhir} \\ \hline
Nilai \textit{objective function} akhir terbaik & #4 \\ \hline
Grafik objektif maksimum dan rata-rata terhadap iterasi & \hasil{nomor gambar/rujukan grafik} \\ \hline
Jumlah populasi & #5 \\ \hline
Banyak iterasi & #6 \\ \hline
Durasi pencarian & #7 detik \\ \hline
\akhirtabel
}
% Argumen: test case, label konfigurasi, nomor percobaan, awalan berkas gambar, nilai akhir, populasi, iterasi, durasi
\newcommand{\percobaanGAGambar}[8]{%
\subsubsection{Test case #1 --- #2 --- Percobaan #3}
\awaltabel
\textit{State} awal (populasi awal) & \includegraphics[width=\linewidth]{gambar/pengujian/ga/#1/#4_initial_state_3d.png} \\ \hline
\textit{State} akhir terbaik & \includegraphics[width=\linewidth]{gambar/pengujian/ga/#1/#4_final_state_3d.png} \\ \hline
Nilai \textit{objective function} akhir terbaik & #5 \\ \hline
Grafik objektif maksimum dan rata-rata terhadap iterasi & \includegraphics[width=\linewidth]{gambar/pengujian/ga/#1/#4_experiment.png} \\ \hline
Jumlah populasi & #6 \\ \hline
Banyak iterasi & #7 \\ \hline
Durasi pencarian & #8 detik \\ \hline
\akhirtabel
}
\newcommand{\tigaPercobaanHC}[1]{%
\percobaanHC{#1}{1}
\percobaanHC{#1}{2}
\percobaanHC{#1}{3}
}
\newcommand{\tigaPercobaanSideways}[1]{%
\percobaanSideways{#1}{1}
\percobaanSideways{#1}{2}
\percobaanSideways{#1}{3}
}
\newcommand{\tigaPercobaanRestart}[1]{%
\percobaanRestart{#1}{1}
\percobaanRestart{#1}{2}
\percobaanRestart{#1}{3}
}
\newcommand{\tigaPercobaanSA}[1]{%
\percobaanSA{#1}{1}
\percobaanSA{#1}{2}
\percobaanSA{#1}{3}
}
\newcommand{\tigaPercobaanGA}[2]{%
\percobaanGA{#1}{#2}{1}
\percobaanGA{#1}{#2}{2}
\percobaanGA{#1}{#2}{3}
}

\section{Eksperimen Algoritma \textit{Hill-climbing}}

\subsection{\textit{Steepest Ascent Hill-climbing}}
\percobaanSteepest{tc8}{1}{run1_seed42}{98}{0.017746}{2}
\percobaanSteepest{tc8}{2}{run2_seed43}{98}{0.015894}{2}
\percobaanSteepest{tc8}{3}{run3_seed44}{96}{0.014675}{2}
\percobaanSteepest{tc9}{1}{run1_seed42}{101}{0.010037}{1}
\percobaanSteepest{tc9}{2}{run2_seed43}{101}{0.010593}{1}
\percobaanSteepest{tc9}{3}{run3_seed44}{96}{0.007649}{1}
\percobaanSteepest{tc10}{1}{run1_seed42}{82}{0.00552}{1}
\percobaanSteepest{tc10}{2}{run2_seed43}{87}{0.006447}{1}
\percobaanSteepest{tc10}{3}{run3_seed44}{83}{0.012795}{2}


\subsection{\textit{Hill-climbing with Sideways Move}}
\percobaanSidewaysGambar{tc8}{1}{run1_seed42}{103}{0.101828}{12}{10}
\percobaanSidewaysGambar{tc8}{2}{run2_seed43}{98}{0.094503}{11}{10}
\percobaanSidewaysGambar{tc8}{3}{run3_seed44}{96}{0.070576}{11}{10}
\percobaanSidewaysGambar{tc9}{1}{run1_seed42}{101}{0.130282}{10}{10}
\percobaanSidewaysGambar{tc9}{2}{run2_seed43}{101}{0.101959}{11}{10}
\percobaanSidewaysGambar{tc9}{3}{run3_seed44}{101}{0.099453}{11}{10}
\percobaanSidewaysGambar{tc10}{1}{run1_seed42}{82}{0.055903}{10}{10}
\percobaanSidewaysGambar{tc10}{2}{run2_seed43}{87}{0.065628}{10}{10}
\percobaanSidewaysGambar{tc10}{3}{run3_seed44}{83}{0.075649}{11}{10}

\subsection{\textit{Random Restart Hill-climbing}}
\percobaanRestartGambar{tc8}{1}{run1_seed42}{103}{0.090698}{3}{2, 1, 2, 1 (rata-rata: 1,5)}
\percobaanRestartGambar{tc8}{2}{run2_seed43}{98}{0.079565}{3}{2, 1, 1, 1 (rata-rata: 1,25)}
\percobaanRestartGambar{tc8}{3}{run3_seed44}{103}{0.074867}{3}{2, 1, 2, 1 (rata-rata: 1,5)}
\percobaanRestartGambar{tc9}{1}{run1_seed42}{101}{0.057637}{3}{1, 2, 1, 1 (rata-rata: 1,25)}
\percobaanRestartGambar{tc9}{2}{run2_seed43}{101}{0.074722}{3}{1, 2, 1, 2 (rata-rata: 1,5)}
\percobaanRestartGambar{tc9}{3}{run3_seed44}{101}{0.055818}{3}{1, 1, 1, 1 (rata-rata: 1)}
\percobaanRestartGambar{tc10}{1}{run1_seed42}{97}{0.053462}{3}{1, 2, 1, 1 (rata-rata: 1,25)}
\percobaanRestartGambar{tc10}{2}{run2_seed43}{88}{0.055554}{3}{1, 1, 2, 1 (rata-rata: 1,25)}
\percobaanRestartGambar{tc10}{3}{run3_seed44}{83}{0.075717}{3}{2, 2, 1, 2 (rata-rata: 1,75)}

\subsection{\textit{Stochastic Hill-climbing}}
\percobaanStochasticGambar{tc8}{1}{run1_seed42}{98}{2.546581}{10.000}
\percobaanStochasticGambar{tc8}{2}{run2_seed43}{98}{2.677560}{10.000}
\percobaanStochasticGambar{tc8}{3}{run3_seed44}{103}{1.668485}{10.000}
\percobaanStochasticGambar{tc9}{1}{run1_seed42}{101}{11.263783}{10.000}
\percobaanStochasticGambar{tc9}{2}{run2_seed43}{101}{14.746734}{10.000}
\percobaanStochasticGambar{tc9}{3}{run3_seed44}{96}{11.694191}{10.000}
\percobaanStochasticGambar{tc10}{1}{run1_seed42}{82}{3.434454}{10.000}
\percobaanStochasticGambar{tc10}{2}{run2_seed43}{87}{9.042528}{10.000}
\percobaanStochasticGambar{tc10}{3}{run3_seed44}{83}{3.207028}{10.000}

\section{Eksperimen Algoritma \textit{Simulated Annealing}}
\percobaanSAGambar{tc8}{1}{run1_seed42}{103}{2}{0.332440}
\percobaanSAGambar{tc8}{2}{run2_seed43}{103}{2}{0.338796}
\percobaanSAGambar{tc8}{3}{run3_seed44}{103}{4}{0.337923}
\percobaanSAGambar{tc9}{1}{run1_seed42}{103}{2}{1.752310}
\percobaanSAGambar{tc9}{2}{run2_seed43}{110}{5}{1.432849}
\percobaanSAGambar{tc9}{3}{run3_seed44}{103}{4}{1.663256}
\percobaanSAGambar{tc10}{1}{run1_seed42}{87}{3}{0.407715}
\percobaanSAGambar{tc10}{2}{run2_seed43}{87}{4}{0.469470}
\percobaanSAGambar{tc10}{3}{run3_seed44}{84}{3}{0.475188}

\section{Eksperimen Algoritma \textit{Genetic Algorithm}}

Setiap konfigurasi dijalankan tiga kali. Konfigurasi populasi 50 dan iterasi 100 dapat muncul pada kedua rangkaian karena merupakan konfigurasi acuan yang sama; hasil tiga \textit{run} yang sama dapat dirujuk pada kedua rangkaian apabila tidak dijalankan ulang.

\subsection{Populasi Tetap, Variasi Banyak Iterasi}
\percobaanGAGambar{tc8}{Populasi 50, iterasi 50}{1}{pop50_it50_run1_seed42}{103}{50}{50}{0.308149}
\percobaanGAGambar{tc8}{Populasi 50, iterasi 50}{2}{pop50_it50_run2_seed43}{103}{50}{50}{0.302332}
\percobaanGAGambar{tc8}{Populasi 50, iterasi 50}{3}{pop50_it50_run3_seed44}{103}{50}{50}{0.337453}
\percobaanGAGambar{tc8}{Populasi 50, iterasi 100}{1}{pop50_it100_run1_seed42}{103}{50}{100}{0.684268}
\percobaanGAGambar{tc8}{Populasi 50, iterasi 100}{2}{pop50_it100_run2_seed43}{103}{50}{100}{0.656125}
\percobaanGAGambar{tc8}{Populasi 50, iterasi 100}{3}{pop50_it100_run3_seed44}{103}{50}{100}{0.606558}
\percobaanGAGambar{tc8}{Populasi 50, iterasi 200}{1}{pop50_it200_run1_seed42}{103}{50}{200}{1.963922}
\percobaanGAGambar{tc8}{Populasi 50, iterasi 200}{2}{pop50_it200_run2_seed43}{103}{50}{200}{1.486151}
\percobaanGAGambar{tc8}{Populasi 50, iterasi 200}{3}{pop50_it200_run3_seed44}{103}{50}{200}{1.417510}
\percobaanGAGambar{tc9}{Populasi 50, iterasi 50}{1}{pop50_it50_run1_seed42}{114}{50}{50}{0.614781}
\percobaanGAGambar{tc9}{Populasi 50, iterasi 50}{2}{pop50_it50_run2_seed43}{110}{50}{50}{0.660053}
\percobaanGAGambar{tc9}{Populasi 50, iterasi 50}{3}{pop50_it50_run3_seed44}{110}{50}{50}{0.634885}
\percobaanGAGambar{tc9}{Populasi 50, iterasi 100}{1}{pop50_it100_run1_seed42}{114}{50}{100}{1.219158}
\percobaanGAGambar{tc9}{Populasi 50, iterasi 100}{2}{pop50_it100_run2_seed43}{110}{50}{100}{1.094391}
\percobaanGAGambar{tc9}{Populasi 50, iterasi 100}{3}{pop50_it100_run3_seed44}{110}{50}{100}{1.047287}
\percobaanGAGambar{tc9}{Populasi 50, iterasi 200}{1}{pop50_it200_run1_seed42}{114}{50}{200}{1.907098}
\percobaanGAGambar{tc9}{Populasi 50, iterasi 200}{2}{pop50_it200_run2_seed43}{110}{50}{200}{2.001291}
\percobaanGAGambar{tc9}{Populasi 50, iterasi 200}{3}{pop50_it200_run3_seed44}{110}{50}{200}{1.882865}
\percobaanGAGambar{tc10}{Populasi 50, iterasi 50}{1}{pop50_it50_run1_seed42}{97}{50}{50}{0.388107}
\percobaanGAGambar{tc10}{Populasi 50, iterasi 50}{2}{pop50_it50_run2_seed43}{97}{50}{50}{0.441870}
\percobaanGAGambar{tc10}{Populasi 50, iterasi 50}{3}{pop50_it50_run3_seed44}{92}{50}{50}{0.526689}
\percobaanGAGambar{tc10}{Populasi 50, iterasi 100}{1}{pop50_it100_run1_seed42}{97}{50}{100}{0.821721}
\percobaanGAGambar{tc10}{Populasi 50, iterasi 100}{2}{pop50_it100_run2_seed43}{97}{50}{100}{0.857301}
\percobaanGAGambar{tc10}{Populasi 50, iterasi 100}{3}{pop50_it100_run3_seed44}{92}{50}{100}{0.809399}
\percobaanGAGambar{tc10}{Populasi 50, iterasi 200}{1}{pop50_it200_run1_seed42}{97}{50}{200}{1.567163}
\percobaanGAGambar{tc10}{Populasi 50, iterasi 200}{2}{pop50_it200_run2_seed43}{97}{50}{200}{1.627509}
\percobaanGAGambar{tc10}{Populasi 50, iterasi 200}{3}{pop50_it200_run3_seed44}{92}{50}{200}{1.655211}

\subsection{Iterasi Tetap, Variasi Jumlah Populasi}
\percobaanGAGambar{tc8}{Populasi 20, iterasi 100}{1}{pop20_it100_run1_seed42}{103}{20}{100}{0.256599}
\percobaanGAGambar{tc8}{Populasi 20, iterasi 100}{2}{pop20_it100_run2_seed43}{103}{20}{100}{0.228431}
\percobaanGAGambar{tc8}{Populasi 20, iterasi 100}{3}{pop20_it100_run3_seed44}{103}{20}{100}{0.298518}
\percobaanGAGambar{tc8}{Populasi 50, iterasi 100}{1}{pop50_it100_run1_seed42}{103}{50}{100}{0.684268}
\percobaanGAGambar{tc8}{Populasi 50, iterasi 100}{2}{pop50_it100_run2_seed43}{103}{50}{100}{0.656125}
\percobaanGAGambar{tc8}{Populasi 50, iterasi 100}{3}{pop50_it100_run3_seed44}{103}{50}{100}{0.606558}
\percobaanGAGambar{tc8}{Populasi 100, iterasi 100}{1}{pop100_it100_run1_seed42}{103}{100}{100}{1.741821}
\percobaanGAGambar{tc8}{Populasi 100, iterasi 100}{2}{pop100_it100_run2_seed43}{103}{100}{100}{1.776312}
\percobaanGAGambar{tc8}{Populasi 100, iterasi 100}{3}{pop100_it100_run3_seed44}{103}{100}{100}{1.591266}
\percobaanGAGambar{tc9}{Populasi 20, iterasi 100}{1}{pop20_it100_run1_seed42}{110}{20}{100}{0.380015}
\percobaanGAGambar{tc9}{Populasi 20, iterasi 100}{2}{pop20_it100_run2_seed43}{105}{20}{100}{0.369793}
\percobaanGAGambar{tc9}{Populasi 20, iterasi 100}{3}{pop20_it100_run3_seed44}{103}{20}{100}{0.340892}
\percobaanGAGambar{tc9}{Populasi 50, iterasi 100}{1}{pop50_it100_run1_seed42}{114}{50}{100}{1.219158}
\percobaanGAGambar{tc9}{Populasi 50, iterasi 100}{2}{pop50_it100_run2_seed43}{110}{50}{100}{1.094391}
\percobaanGAGambar{tc9}{Populasi 50, iterasi 100}{3}{pop50_it100_run3_seed44}{110}{50}{100}{1.047287}
\percobaanGAGambar{tc9}{Populasi 100, iterasi 100}{1}{pop100_it100_run1_seed42}{114}{100}{100}{2.355519}
\percobaanGAGambar{tc9}{Populasi 100, iterasi 100}{2}{pop100_it100_run2_seed43}{110}{100}{100}{2.311342}
\percobaanGAGambar{tc9}{Populasi 100, iterasi 100}{3}{pop100_it100_run3_seed44}{110}{100}{100}{2.464745}
\percobaanGAGambar{tc10}{Populasi 20, iterasi 100}{1}{pop20_it100_run1_seed42}{92}{20}{100}{0.258467}
\percobaanGAGambar{tc10}{Populasi 20, iterasi 100}{2}{pop20_it100_run2_seed43}{92}{20}{100}{0.353925}
\percobaanGAGambar{tc10}{Populasi 20, iterasi 100}{3}{pop20_it100_run3_seed44}{92}{20}{100}{0.361047}
\percobaanGAGambar{tc10}{Populasi 50, iterasi 100}{1}{pop50_it100_run1_seed42}{97}{50}{100}{0.821721}
\percobaanGAGambar{tc10}{Populasi 50, iterasi 100}{2}{pop50_it100_run2_seed43}{97}{50}{100}{0.857301}
\percobaanGAGambar{tc10}{Populasi 50, iterasi 100}{3}{pop50_it100_run3_seed44}{92}{50}{100}{0.809399}
\percobaanGAGambar{tc10}{Populasi 100, iterasi 100}{1}{pop100_it100_run1_seed42}{97}{100}{100}{1.995256}
\percobaanGAGambar{tc10}{Populasi 100, iterasi 100}{2}{pop100_it100_run2_seed43}{97}{100}{100}{1.818441}
\percobaanGAGambar{tc10}{Populasi 100, iterasi 100}{3}{pop100_it100_run3_seed44}{97}{100}{100}{1.809598}

\section{Analisis Hasil Eksperimen}

Analisis hasil eksperimen dapat disusun dengan membandingkan kedekatan setiap algoritma terhadap nilai terbaik yang diperoleh, konsistensi nilai akhir antar-\textit{run}, dan durasi pencarian. Untuk \textit{Genetic Algorithm}, analisis juga membahas pengaruh variasi jumlah iterasi serta jumlah populasi terhadap nilai akhir dan durasi. Jelaskan pula keterkaitan perilaku tersebut dengan mekanisme eksplorasi dan eksploitasi masing-masing algoritma.
